\documentclass[12font]{article}
\title{\textbf{ Pseudo Flat Band }}
\author{Chen-Chih Wang  }
\date{\textit{Department of physics, National Tsing Hua University, Hsinchu 30013, Taiwan} \\ (Dated: \today)}
\usepackage[margin=2cm]{geometry}
%\usepackage[fleqn]{amsmath}
%\usepackage{CJK,amsmath}
\usepackage{amsfonts,amssymb}
\usepackage{fancyhdr}
%\pagestyle{fancy}
\usepackage{textcomp}
\usepackage{graphicx}
\usepackage{float}
\usepackage{color}
\usepackage{tikz}
\usepackage{multicol}
\usepackage{threeparttable}
\usepackage{amsthm}
\usepackage{mathtools}
\usepackage{hyperref}
\urlstyle{same}
\usepackage{caption}
\usepackage{subcaption}
\usepackage{dsfont}

\newcommand{\bra}[1]{\left\langle #1 \right|}
\newcommand{\ket}[1]{\left| #1 \right\rangle}
\newcommand{\anglelr}[1]{\left\langle #1 \right\rangle}
\newcommand{\braket}[2]{\left\langle #1 \right|\left. #2 \right\rangle}
\newcommand{\dif}[1]{\frac{\mathrm{d}}{\mathrm{d}#1}}
\newcommand{\diff}[2]{\frac{\mathrm{d}#1}{\mathrm{d}#2}}
\newcommand{\gd}{\boldsymbol{\nabla}}
\newcommand{\dive}{\boldsymbol{\nabla}\cdot}
\newcommand{\curl}{\boldsymbol{\nabla}\times}
\newcommand{\lr}[1]{\left(  #1  \right)}
\newcommand{\lrm}[1]{\left[  #1  \right]}
\newcommand{\lrg}[1]{\left\{  #1  \right\}}
\newcommand{\abs}[1]{\left|  #1  \right|}
\newcommand{\de}{\partial}
\newcommand{\pdif}[1]{\frac{\partial}{\partial#1}}
\newcommand{\pdiff}[2]{\frac{\partial#1}{\partial#2}}
\newcommand{\mrm}[1]{\mathrm{#1}}
\newcommand{\bo}[1]{\mathbf{#1}}
\newcommand{\bog}[1]{\boldsymbol{#1}}
\newcommand{\vint}[3]{\int_{#1}#2\cdot\mathrm{d}\mathbf{#3}}
\newcommand{\voint}[3]{\oint_{#1}#2\cdot\mathrm{d}\mathbf{#3}}
\newcommand{\Vint}{\int\mrm{d}^3\bo{r}'}
\newcommand{\rr}{\mathbf{r} - \mathbf{r}'}
\newcommand{\arr}{\left|\mathbf{r} - \mathbf{r}'\right|}

\begin{document}
\maketitle

%\fontsize{12pt}{18pt}\selectfont


\begin{abstract}
  I found that one can obtain a flat band solution by pure mathematical trick. In other words, for some systems there seems to be two different solution to the band structure. If one can fund a local particle representation to diagonalize the Hamiltonian, then one can find the flat band, and which is dubbed the \emph{pseudo flat band} in this note.
\end{abstract}

\section{Ising model}
Starting with the one dimensional Ising Hamiltonian without the external magnetic field,
\[
\hat{H}_{\mrm{Ising}} = -\sum_{i}S^x_i S^x_{i+1}
\]
with the one dimensional Jordan-Wigner (JW) transformation,
\[
\left\{
\begin{aligned}
& S^x_i = \lrm{\prod_{j<i}\lr{1-2c_j^\dagger c_j}}\lr{c_i^\dagger + c_i} \\
& S^y_i = \lrm{\prod_{j<i}\lr{1-2c_j^\dagger c_j}}i\lr{c_i^\dagger - c_i} \\
& S^z_i = 1 - 2c_i^\dagger c_i
\end{aligned}
\right.
\]
The Ising Hamiltonian can be written as,
\[
\hat{H}_{\mrm{Ising}} = -\sum_{i}\lr{c_i^\dagger - c_i}\lr{c_{i+1}^\dagger + c_{i+1}} \equiv -\sum_{i}i\alpha_{i+1}\beta_{i}
\]
where 
\[
\alpha_i \equiv c_i^\dagger + c_i \;,\; \beta_i \equiv i\lr{c_i^\dagger - c_i}
\]
Similar to the Kitaev's trick, we can define an Ising operator
\[\hat{u}_i \equiv i\alpha_{i+1}\beta_i\]
we can check that $\hat{u}_i^2 = \mathds{1}$. Moreover, each $\hat{u}_i$ commutes with all the other one, that is $[\hat{u}_i,\hat{u}_j] = 0$. Thus each $\hat{u}_i$ commutes with the Hamiltonian and hence its eigenvalue is a good quantum number. Because of these properties, such an abelian $\mathbb{Z}_2$ operators can be represented by fermion operators $\eta$ in such a way that,
\[
\hat{u}_i \equiv 1 - 2\eta_i^\dagger \eta_i
\]
we can check that this representation preserves the ablian as well as the $\mathbb{Z}_2$ properties. Plugging this representation into the original Hamiltonian gives,
\[
\hat{H}_{\mrm{Ising}} = -\sum_{i}\hat{u}_i = \sum_{i}\lr{2\eta_i^\dagger\eta_i - 1} = \sum_{k}2\eta_k^\dagger\eta_k - N
\]
in the final step we Fourier transform the fermion, which is redundant in some sense since the Hamiltonian is already diagonal before doing that. Eventually we obtain the flat band $\epsilon(k) = 2$, which is a constant. That makes sense since the fermion $\eta$, the elementary excitation in this diagonalization scheme, is completely local in the real space, i.e. no hopping term can make them move to other lattice sites. 

On the other hand, if we diagonalize the Hamiltonian in the JW fermion representation, we have,
\begin{align*}
  \hat{H}_{\mrm{Ising}} &= \sum_{i}\lrm{c_{i+1}^\dagger c_i^\dagger - c_{i+1}^\dagger c_i + \mrm{H.c.}} \\
  &= \sum_{k}\lrm{e^{-ik}c_{k}^\dagger c_{-k}^\dagger - e^{-ik} c_k^\dagger c_k + \mrm{H.c.}} \\
  &= \sum_{k}\lrm{e^{-ik}c_{k}^\dagger c_{-k}^\dagger - \cos k c_k^\dagger c_k + \mrm{H.c.}} \\
  &= \begin{pmatrix}
       c_k^\dagger & c_{-k} 
     \end{pmatrix}\begin{pmatrix}
                    -\cos k & e^{-ik} \\
                    e^{ik} & \cos k 
                  \end{pmatrix}\begin{pmatrix}
                                 c_k \\
                                 c_{-k}^\dagger 
                               \end{pmatrix} + E_0
\end{align*}
The matrix sandwiched between spinors in the last line can be further represented as,
\[
\begin{pmatrix}
-\cos k & e^{-ik} \\
e^{ik} & \cos k 
\end{pmatrix} = \cos k\sigma^x + \sin k \sigma^y - \cos k \sigma^z
\]
thus the dispersion is
\begin{equation}\label{real band structure.e}
  \epsilon_{\pm}\lr{k} = \pm\sqrt{1 + \cos^2 k}
\end{equation}
which is obviously dispersive and doesn't consist with the flat band result solved by $\eta$ fermion representation.


\begin{thebibliography}{99}

    \bibitem{}{}

\end{thebibliography}

\end{document}
